\scanpage{174}
\mspage{32}

\inserted{en degrés $>0$} est nulle, donc $\xi=0$, \unsure{qed.|cqfd.}

% slanted
\reversemarginpar\margin{C'est déjà dans SGA4 I 8}\normalmarginpar

\uline{Proposition} 2.15. Soit $X$ un objet de Cat. Conditions équivalentes :

a) $X$ est filtrante, i.e. $X$ est $0$-connexe et $X^{\unsure{o|0}}$ satisfait les conditions \unsure{équiv.} du th. 2.6 ($X^{o\,\wedge}$ est loc. tot. $0$ connexe).

b) \deleted{\unsure{tous les objets foncteurs}} L'objet final de $X^\wedge$ est ind-représentable.

\uline{Dém.} On sait qu'on a

\[
e_{X^\wedge}\simeq\varinjlim_{x\in X}x,\tag*{(2.15.1)}
\]

limite inductive prise dans $X^\wedge$. En effet, cela signifie aussi que pour tout objet $a$ de $X$, on a

\[
\varinjlim_{x\in X}\underbrace{x(a)}_{\operatorname{Hom}_{X}(a,\unsure{X|x})}\xrightarrow{\sim}e
\quad\text{(\unsure{ens. un elt|ens. ensbl} ponctuel)}.
\]

Mais \unsure{considérons} l'\deleted{\unsure{espace}} \inserted{objet} \deleted{étale} $X'={}_{a}\backslash X\deleted{\unsure{_{/a}}}\unsure{\mathbin{-}|\longrightarrow}X$ coétale sur $X$, qui est \inserted{au \unsure{Cat-|cat-}} \unsure{cofibrée|cofibré} sur $X$, on sait que l'on a

\[
\pi_0(X'\deleted{\unsure{_{/x}}})\simeq\varinjlim_{x\in X}\deleted{\unsure{X_x}}\pi_0(X'_x),\quad\text{or}
\]

$X'_x=\operatorname{Hom}(a,x)$, donc on trouve

\[
\varinjlim_{x\in X}\operatorname{Hom}(a,x)=\pi_0({}_{a}\backslash X).=e
\]

car $\,{}_{a}\backslash X$ est $0$-connexe

La formule (2.15.1.) montre que si

% archive sheet number: 174
